\section{Scheme theoretic image}
\label{section-scheme-theoretic-image}

\noindent
Caution: Some of the material in this section is ultra-general and
behaves differently from what you might expect.

\begin{lemma}
\label{lemma-scheme-theoretic-image}
Let $f : X \to Y$ be a morphism of schemes. There exists a closed
subscheme $Z \subset Y$ such that $f$ factors through $Z$ and such
that for any other closed subscheme $Z' \subset Y$ such that $f$
factors through $Z'$ we have $Z \subset Z'$.
\end{lemma}

\begin{proof}
Let $\mathcal{I} = \Ker(\mathcal{O}_Y \to f_*\mathcal{O}_X)$.
If $\mathcal{I}$ is quasi-coherent then we just take $Z$ to be the
closed subscheme determined by $\mathcal{I}$, see
Lemma \ref{lemma-closed-immersion-bijection-ideals}. This works by
Lemma \ref{lemma-closed-immersion-ideals}.
In general the same lemma requires us to show that there exists
a largest quasi-coherent sheaf of ideals $\mathcal{I}'$ contained in
$\mathcal{I}$.
This follows from Lemma \ref{lemma-largest-quasi-coherent-subsheaf}.
\end{proof}

\begin{definition}
\label{definition-scheme-theoretic-image}
Let $f : X \to Y$ be a morphism of schemes. The {\it scheme theoretic image}
of $f$ is the smallest closed subscheme $Z \subset Y$ through which $f$
factors, see Lemma \ref{lemma-scheme-theoretic-image} above.
\end{definition}

\noindent
For a morphism $f : X \to Y$ of schemes with scheme theoretic image $Z$
we often denote $f : X \to Z$ the factorization of $f$
through its scheme theoretic image. If the morphism $f$ is not
quasi-compact, then (in general)
\begin{enumerate}
\item the set theoretic inclusion $\overline{f(X)} \subset Z$
is not an equality, i.e., $f(X) \subset Z$ is not a dense subset, and
\item the construction of the scheme theoretic image does not commute with
restriction to open subschemes to $Y$.
\end{enumerate}
In Examples, Section \ref{examples-section-scheme-theoretic-image}
the reader finds an example for both phenomena.
These phenomena can arise even for immersions, see
Examples, Section \ref{examples-section-strange-immersion}.
However, the next lemma shows that both disasters are avoided
when the morphism is quasi-compact.

\begin{lemma}
\label{lemma-quasi-compact-scheme-theoretic-image}
Let $f : X \to Y$ be a morphism of schemes.
Let $Z \subset Y$ be the scheme theoretic image of $f$.
If $f$ is quasi-compact then
\begin{enumerate}
\item the sheaf of ideals
$\mathcal{I} = \Ker(\mathcal{O}_Y \to f_*\mathcal{O}_X)$
is quasi-coherent,
\item the scheme theoretic image $Z$ is the closed subscheme
determined by $\mathcal{I}$,
\item for any open $U \subset Y$ the scheme theoretic image of
$f|_{f^{-1}(U)} : f^{-1}(U) \to U$ is equal to $Z \cap U$, and
\item the image $f(X) \subset Z$ is a dense subset of $Z$, in other
words the morphism $X \to Z$ is dominant
(see Definition \ref{definition-dominant}).
\end{enumerate}
\end{lemma}

\begin{proof}
Part (4) follows from part (3). To show (3) it suffices
to prove (1) since the formation of $\mathcal{I}$ commutes with restriction to
open subschemes of $Y$. And if (1) holds then in the proof of
Lemma \ref{lemma-scheme-theoretic-image}
we showed (2). Thus it suffices to prove that $\mathcal{I}$ is quasi-coherent.
Since the property of being quasi-coherent is
local we may assume $Y$ is affine. As $f$ is quasi-compact,
we can find a finite affine open covering
$X = \bigcup_{i = 1, \ldots, n} U_i$. Denote $f'$ the composition
$$
X' = \coprod U_i \longrightarrow X \longrightarrow Y.
$$
Then $f_*\mathcal{O}_X$ is a subsheaf of $f'_*\mathcal{O}_{X'}$,
and hence $\mathcal{I} = \Ker(\mathcal{O}_Y \to f'_*\mathcal{O}_{X'})$.
By Schemes, Lemma \ref{schemes-lemma-push-forward-quasi-coherent}
the sheaf $f'_*\mathcal{O}_{X'}$ is quasi-coherent on $Y$. Hence we win.
\end{proof}

\begin{example}
\label{example-scheme-theoretic-image}
If $A \to B$ is a ring map with kernel $I$, then the scheme theoretic image
of $\Spec(B) \to \Spec(A)$ is the closed subscheme
$\Spec(A/I)$ of $\Spec(A)$. This follows from
Lemma \ref{lemma-quasi-compact-scheme-theoretic-image}.
\end{example}

\noindent
If the morphism is quasi-compact, then the scheme theoretic image only
adds points which are specializations of points in the image.

\begin{lemma}
\label{lemma-reach-points-scheme-theoretic-image}
Let $f : X \to Y$ be a quasi-compact morphism.
Let $Z$ be the scheme theoretic image of $f$.
Let $z \in Z$\footnote{By
Lemma \ref{lemma-quasi-compact-scheme-theoretic-image} set-theoretically
$Z$ agrees with the closure of $f(X)$ in $Y$.}.
There exists a valuation ring $A$ with
fraction field $K$ and a commutative diagram
$$
\xymatrix{
\Spec(K) \ar[rr] \ar[d] & & X \ar[d] \ar[ld] \\
\Spec(A) \ar[r] & Z \ar[r] & Y
}
$$
such that the closed point of $\Spec(A)$ maps to $z$. In particular
any point of $Z$ is the specialization of a point of $f(X)$.
\end{lemma}

\begin{proof}
Let $z \in \Spec(R) = V \subset Y$ be an affine open
neighbourhood of $z$. By
Lemma \ref{lemma-quasi-compact-scheme-theoretic-image}
the intersection $Z \cap V$ is the scheme theoretic image of
$f^{-1}(V) \to V$. Hence we may replace $Y$ by $V$
and assume $Y = \Spec(R)$ is affine.
In this case $X$ is quasi-compact as $f$ is quasi-compact.
Say $X = U_1 \cup \ldots \cup U_n$
is a finite affine open covering. Write $U_i = \Spec(A_i)$.
Let $I = \Ker(R \to A_1 \times \ldots \times A_n)$.
By Lemma \ref{lemma-quasi-compact-scheme-theoretic-image}
again we see that $Z$ corresponds to the closed subscheme
$\Spec(R/I)$ of $Y$. If $\mathfrak p \subset R$ is
the prime corresponding to $z$, then we see that
$I_{\mathfrak p} \subset R_{\mathfrak p}$ is not an
equality. Hence (as localization is exact, see
Algebra, Proposition \ref{algebra-proposition-localization-exact})
we see that
$R_{\mathfrak p} \to
(A_1)_{\mathfrak p} \times \ldots \times (A_n)_{\mathfrak p}$
is not zero. Hence one of the rings $(A_i)_{\mathfrak p}$ is not zero.
Hence there exists an $i$ and a prime $\mathfrak q_i \subset A_i$
lying over a prime $\mathfrak p_i \subset \mathfrak p$.
By Algebra, Lemma \ref{algebra-lemma-dominate} we can choose a valuation ring
$A \subset K = \kappa(\mathfrak q_i)$ dominating
the local ring
$R_{\mathfrak p}/\mathfrak p_iR_{\mathfrak p} \subset \kappa(\mathfrak q_i)$.
This gives the desired diagram. Some details omitted.
\end{proof}

\begin{lemma}
\label{lemma-factor-factor}
Let
$$
\xymatrix{
X_1 \ar[d] \ar[r]_{f_1} & Y_1 \ar[d] \\
X_2 \ar[r]^{f_2} & Y_2
}
$$
be a commutative diagram of schemes. Let $Z_i \subset Y_i$, $i = 1, 2$ be
the scheme theoretic image of $f_i$. Then the morphism
$Y_1 \to Y_2$ induces a morphism $Z_1 \to Z_2$ and a
commutative diagram
$$
\xymatrix{
X_1 \ar[r] \ar[d] & Z_1 \ar[d] \ar[r] & Y_1 \ar[d] \\
X_2 \ar[r] & Z_2 \ar[r] & Y_2
}
$$
\end{lemma}

\begin{proof}
The scheme theoretic inverse image of $Z_2$ in $Y_1$
is a closed subscheme of $Y_1$ through
which $f_1$ factors. Hence $Z_1$ is contained in this.
This proves the lemma.
\end{proof}

\begin{lemma}
\label{lemma-scheme-theoretic-image-reduced}
Let $f : X \to Y$ be a morphism of schemes.
If $X$ is reduced, then the scheme theoretic image of $f$ is
the reduced induced scheme structure on $\overline{f(X)}$.
\end{lemma}

\begin{proof}
This is true because the reduced induced scheme structure on $\overline{f(X)}$
is clearly the smallest closed subscheme of $Y$ through which $f$ factors,
see
Schemes, Lemma \ref{schemes-lemma-map-into-reduction}.
\end{proof}

\begin{lemma}
\label{lemma-scheme-theoretic-image-of-partial-section}
Let $f : X \to Y$ be a separated morphism of schemes.
Let $V \subset Y$ be a retrocompact open. Let $s : V \to X$
be a morphism such that $f \circ s = \text{id}_V$.
Let $Y'$ be the scheme theoretic image of $s$.
Then $Y' \to Y$ is an isomorphism over $V$.
\end{lemma}

\begin{proof}
The assumption that $V$ is retrocompact in $Y$
(Topology, Definition \ref{topology-definition-quasi-compact})
means that $V \to Y$ is a quasi-compact morphism.
By Schemes, Lemma \ref{schemes-lemma-quasi-compact-permanence}
the morphism $s : V \to X$ is quasi-compact.
Hence the construction of the scheme theoretic image $Y'$
of $s$ commutes with restriction to opens by
Lemma \ref{lemma-quasi-compact-scheme-theoretic-image}.
In particular, we see that $Y' \cap f^{-1}(V)$ is the
scheme theoretic image of a section of the separated
morphism $f^{-1}(V) \to V$. Since a section of a separated
morphism is a closed immersion
(Schemes, Lemma \ref{schemes-lemma-section-immersion}),
we conclude that
$Y' \cap f^{-1}(V) \to V$ is an isomorphism as desired.
\end{proof}







